\section{Einleitung}\label{sec:einleitung}

\subsection{Das Problem: stille Übergänge}

Ein Sprachmodell kann in einer Sitzung mehr Seiten mathematischen Text erzeugen, als ein Mensch an einem Tag sorgfältig lesen
kann. Die Fehler, die in einem solchen Bestand gefährlich werden, sind dabei meist nicht die offensichtlichen. Eine falsch
ausmultiplizierte Klammer fällt bei der nächsten Rechnung auf. Schwerer zu sehen sind Übergänge, die im Fließtext keine Spur
hinterlassen:
\begin{itemize}[nosep]
\item \emph{Beförderung:} Eine Aussage, die nur unter einer Bedingung gilt oder nur numerisch gezeigt wurde, wird drei Seiten
später als Satz zitiert und als Grundlage weiterverwendet.
\item \emph{Eingabe als Ergebnis:} Ein Wert, der aus der Literatur eingesetzt oder angepasst wurde, erscheint später als
Vorhersage.
\item \emph{Frame-Vermischung:} Zwei Größen werden verglichen oder verrechnet, obwohl sie zu verschiedenen Beschreibungsrahmen
gehören und keine Abbildung zwischen ihnen erklärt ist.
\end{itemize}
Keiner dieser Schritte ist für sich ein Rechenfehler, und alle drei sind für ein Sprachmodell naheliegend. Ein Modell setzt
einen Text so fort, dass er zusammenhängend wirkt. Eine Bedingung, die vor zwanzig Absätzen genannt wurde, trägt wenig zum
Zusammenhang des nächsten Satzes bei; die Fortsetzung glättet sie weg. Das ist kein seltenes Versagen, sondern die erwartbare
Folge eines Verfahrens, das auf flüssige Fortsetzung optimiert ist. Wer mit Sprachmodellen Theorie entwickelt, braucht deshalb
eine Struktur außerhalb des Texts, die festhält, worauf jede Aussage ruht, und die sich nicht durch Umformulieren ändern lässt.

Dieser Artikel beschreibt eine solche Struktur, QAM. Wir stellen sie zuerst als eigenständiges, von jedem Anwendungsgebiet
unabhängiges Modell vor, dann ihre Umsetzung in \LaTeX, Lean und einer Ausführungsmaschine, und erst danach ihren Einsatz in einem
konkreten Forschungsprojekt. Zwei Grundpfeiler der Informatik liefern die Konstruktionsprinzipien. Wir verwenden sie nicht als
historische Einkleidung, sondern weil sich aus ihnen Entwurfsentscheidungen ableiten, die im System tatsächlich umgesetzt sind.

\subsection{Zuse: Rechenwerk und Plan}

Konrad Zuses Z3 (1941) trennte das Rechenwerk, das binäre Gleitkommaoperationen ausführte, vom Speicher und von einem Programm, das
von gelochtem Filmstreifen gelesen wurde~\cite{rojas}. Die Maschine führte aus, was der Streifen vorgab; eine bedingte Verzweigung
besaß sie nicht. Im Plankalkül (entworfen 1942--1945, veröffentlicht 1972~\cite{zuse72,bauer}) ging Zuse den umgekehrten Weg und
beschrieb Rechenpläne unabhängig von einer Maschine: mit strukturierten, aus Bits aufgebauten Datentypen und mit getrennten
Variablenklassen für Eingangswerte ($V$), Zwischenwerte ($Z$) und Resultatwerte ($R$)~\cite{bauer}.

Für die Arbeit mit Sprachmodellen übernehmen wir aus beidem eine Rollenteilung:
\begin{center}
\begin{tabular}{@{}p{0.22\textwidth}p{0.33\textwidth}p{0.36\textwidth}@{}}
\toprule
Rolle & bei Zuse & in einem QAM-Arbeitsablauf\\
\midrule
Ausführung & Rechenwerk der Z3 & Sprachmodell, Computeralgebra, Simulator\\
Plan & Rechenplan im Plankalkül & QAM-Akte und MML-Auftrag: Typen, Frames, Status, Prämissen\\
Eingang vs.\ Resultat & $V$- und $R$-Variablen & Eingabe (\texttt{input}) vs.\ Herleitung (\texttt{derive}); Prüfung \texttt{\textbackslash qamindependent}\\
\bottomrule
\end{tabular}
\end{center}
Das Sprachmodell darf also formulieren, umformen und rechnen. Welche Objekte existieren, welchen Status sie haben und wovon sie
abhängen, legt es nicht durch seinen Text fest, sondern durch Einträge in einer Akte, die eigene Prüfer lesen. Die Analogie hat
eine Grenze, die wir ausdrücklich nennen: Zuse trennte Plan und Ausführung, um Rechnungen mechanisch ausführbar zu machen; wir
trennen sie, um die Ausführung durch einen unzuverlässigen Ausführenden kontrollierbar zu machen.

Im Repository ist diese Teilung auch wörtlich umgesetzt (Abschnitt~\ref{sec:maschine}): Die MML-Anweisung \texttt{RUN\_TAPE}
führt eine binäre Bandmaschine aus, deren Übergangstabelle als QAM-Objekt der Art \texttt{operator} und deren Startkonfiguration als
\texttt{state} deklariert ist. Zwei SystemVerilog-Steuerungen führen denselben Plan aus und werden an
209 Referenzfällen gegen die Python-Ausführung verglichen. Der Plan ist das Objekt, die Ausführungen sind austauschbar. Dass gerade eine Bandmaschine gewählt ist, hat einen inhaltlichen
Grund: Sie ist das ausführbare Gegenstück zu der offenen Frage, ob die mathematische Struktur des Anwendungsprojekts selbst eine
universelle Maschine trägt (Abschnitt~\ref{sec:traeger}).

\subsection{Turing: die Entscheidbarkeitsgrenze}

Turing zeigte 1936, dass es kein mechanisches Verfahren gibt, das für jede Aussage der Prädikatenlogik entscheidet, ob sie
beweisbar ist~\cite{turing}. Rice verallgemeinerte später: Jede nichttriviale semantische Eigenschaft von Programmen ist
unentscheidbar~\cite{rice}. Für ein Prüfsystem, das Forschungstext begleitet, folgt daraus eine nüchterne Entwurfsregel. Es wird
Aussagen geben, die das System weder bestätigen noch widerlegen kann, und zwar nicht nur aus Mangel an Rechenzeit, sondern
grundsätzlich. Ein System, das nur die Ausgänge \glqq gilt\grqq{} und \glqq gilt nicht\grqq{} kennt, muss solche Aussagen einem
der beiden zuschlagen. Beides ist falsch: Das erste ist die Beförderung aus der Fehlerliste oben, das zweite verwirft Wissen,
das unter Bedingungen sehr wohl gilt.

QAM führt deshalb den dritten Ausgang als erstklassigen Bestandteil. Auf der Ebene der Akte ist er die Statusordnung
\(\st{open}<\st{reading}<\st{conditional}<\st{computed}<\st{proved}\) zusammen mit der Regel, dass kein Schluss stärker ist als seine
schwächste Prämisse. Auf der Ebene der Maschine endet jeder Schritt mit \textsc{cont}, \textsc{reject} oder \textsc{guard}, und ein
\textsc{guard} wird mit seinen offenen Bedingungen gespeichert.

\subsection{Beiträge und Aufbau}

Abschnitt~\ref{sec:modell} definiert das Modell: registrierte Aussage, vier Wohlgeformtheitsregeln und offene Bedingungen, die an jeden
Nachfolger vererbt werden. Abschnitt~\ref{sec:latex} beschreibt die \LaTeX-Anbindung, Abschnitt~\ref{sec:guards} die Prüfer und zwei
getrennte Guard-Ebenen mit genau benannter Prüftiefe, Abschnitt~\ref{sec:maschine} Datenbank, Ausführungsmaschine und die Anbindung an
Sprachmodelle über das Model Context Protocol. Erst Abschnitt~\ref{sec:fallstudie} zeigt den Einsatz im Projekt \emph{Quantum Sphaera}
(rund 1450 Seiten Quelltext), mit einem konstruierten und einem protokollierten Beispiel und mit den Fehlern, die QAM \emph{nicht} fing.
Abschnitt~\ref{sec:abgrenzung} grenzt QAM gegen Beweisassistenten ab.
Durchgängig trennen wir drei Ebenen: was formal festgelegt und geprüft ist, was sich im Einsatz bewährt hat, und was QAM bewusst
nicht leistet.
